\begingroup
\color{borrador2}
\textcolor{borrador2}{La siguiente figura separa como alternativa el procesamiento fuera de línea del recorrido de cada consulta, para evitar mezclar fases de construcción y de ejecución.}

\begin{figure}[H]
  \centering
  \begin{tikzpicture}[
    nodo/.style={draw=borrador2, thick, rounded corners=3pt, fill=white,
                 text=borrador2, align=center, font=\footnotesize,
                 minimum width=2.55cm, minimum height=1.05cm},
    dato/.style={nodo, fill=borrador2!8, font=\footnotesize\ttfamily},
    grupo/.style={draw=borrador2, densely dashed, rounded corners=5pt,
                  inner sep=9pt},
    titulo/.style={font=\small\bfseries, text=borrador2},
    flecha/.style={-{Latex[length=2.2mm]}, thick, draw=borrador2},
    doble/.style={{Latex[length=2.2mm]}-{Latex[length=2.2mm]}, thick, draw=borrador2},
  ]
    \node[nodo] (web) at (-6.0,2.1) {Web oficial\\de la EPSJ};
    \node[nodo] (spider) at (-3.0,2.1) {\textit{Spider} Scrapy};
    \node[dato] (grados) at (0,2.1) {grados.json};
    \node[nodo] (chunker) at (3.0,2.1) {Fragmentación y\\deduplicación};
    \node[dato] (chunks) at (6.0,2.1) {chunks.json};

    \node[nodo] (indexer) at (3.0,0) {Incrustación e\\indexación};
    \node[nodo, cylinder, shape border rotate=90, aspect=0.18] (lance) at (6.0,0) {LanceDB};

    \draw[flecha] (web) -- (spider);
    \draw[flecha] (spider) -- (grados);
    \draw[flecha] (grados) -- (chunker);
    \draw[flecha] (chunker) -- (chunks);
    \draw[flecha] (chunks) -- (indexer);
    \draw[flecha] (indexer) -- (lance);

    \node[nodo] (navegador) at (-5.2,-3.2) {Navegador\\{\scriptsize sin cuenta de usuario}};
    \node[nodo, minimum width=5.8cm] (servidor) at (0,-3.2)
      {Servidor web\\{\scriptsize conversación y ámbito}\\{\scriptsize recuperación, generación y verificación}};
    \node[nodo] (ollama) at (5.2,-3.2) {Servidor local\\de inferencia};
    \node[dato, minimum width=3.7cm] (registro) at (5.2,-5.0) {registro\_chat.jsonl};

    \draw[doble] (navegador) -- node[above, font=\scriptsize, text=borrador2]{flujo de eventos} (servidor);
    \draw[doble] (servidor) -- (ollama);
    \draw[doble] (servidor.north east) -- node[above, sloped, font=\scriptsize, text=borrador2]{búsqueda} (lance.south west);
    \draw[flecha] (servidor.south east) -- node[right, font=\scriptsize, text=borrador2]{traza} (registro.west);

    \begin{scope}[on background layer]
      \node[grupo, fit=(web)(spider)(grados)(chunker)(chunks)(indexer)(lance),
            label={[titulo]north:Construcción reproducible del corpus y del índice}] {};
      \node[grupo, fit=(navegador)(servidor)(ollama)(registro),
            label={[titulo]south:Ejecución de una consulta}] {};
    \end{scope}
  \end{tikzpicture}
  \caption[Alternativa de la arquitectura general]{\textcolor{borrador2}{Versión alternativa de la arquitectura general. La parte superior se ejecuta para construir el índice; la inferior se ejecuta al atender cada consulta. La ausencia de cuenta de usuario no implica ausencia de registro técnico. Fuente: Elaboración propia.}}
  \label{fig:arquitectura_pipeline_alt}
\end{figure}
\endgroup
